\chapter{Generation and Citation Validation}
\label{ch:generation}

\section{The prompt}

The answering prompt has two parts: a system instruction holding the rules, and a user message holding the numbered passages and the question \ev{src/vidyarag/generate/prompts.py}. This is the rules text as the model receives it:

\begin{lstlisting}[language=plain, frame=single, backgroundcolor=\color{blue!3}]
You are a study assistant that answers questions strictly from supplied textbook passages.
Rules:
1. Use ONLY the numbered passages provided. Do not add facts from memory, even if you are
   confident they are correct.
2. Cite every factual claim with the marker of the passage supporting it, like [1] or [3].
3. If the passages do not contain the answer, say exactly what is missing rather than
   guessing. A partial answer that names its gap is more useful than a complete-sounding
   one that invents the rest.
4. Never cite a number that was not provided to you.
5. Answer in clear prose for a student revising the topic. Be direct.
\end{lstlisting}

Each passage appears as \code{[n] <citation>} on one line, followed by its text. Every rule has a job. \textbf{Rule~1} makes this RAG rather than a chatbot; ``even if you are confident'' closes the loophole of a true but uncitable remembered fact. \textbf{Rule~2} uses small numbers rather than chunk ids because models ``mangle'' identifiers. \textbf{Rule~3} is the one that matters most: \emph{it makes the generator refuse on its own}. Chapter~\ref{ch:evaluation} shows that it refused every unanswerable question in every profile. \textbf{Rule~4} is a defence in words; the validator below is the defence in code. The prompt is versioned (\code{answer-v1}), and the version enters the trace, the run files and the answer-cache key: ``editing in place silently invalidates every number measured with the previous wording.''

\section{Calling the model}

\code{generate_answer} \ev{src/vidyarag/generate/answer.py} does five things, in order:
\begin{enumerate}
  \item \textbf{No context, no call.} If retrieval returned nothing, the model is never asked (``asking a model to answer with no context is asking it to hallucinate''), and a fixed ``found nothing'' message is returned.
  \item \textbf{Build the prompt and call Gemini} at temperature 0, with \code{max_output_tokens} 2,048. Temperature 0 because ``a benchmark that moves between runs cannot support a claim about a change'', although it is \emph{not} deterministic in practice.
  \item \textbf{Record usage}, through the same helper that grading and decomposition now use.
  \item \textbf{Extract text robustly.} The SDK's \code{response.text} is \code{None} on a truncated response, so the parts are walked directly, rather than letting a truncated answer ``surface as the string \enquote{None}''.
  \item \textbf{Validate citations} (below).
\end{enumerate}
In the committed \code{baseline} run, a generation call averaged 2,400 input and 59 output tokens.

\paragraph{What ``grounded'' means.} \code{GeneratedAnswer.grounded} is \code{bool(citations)}: at least one citation resolved. It says nothing about whether the cited passages support the claims. Checking that is the grader's job, and it happens only in the self-check profiles. There is no \code{try} around the model call, so a rate-limit or network error propagates to the interface. The API now turns it into a 429 or 503, and the demo into a plain message (Chapter~\ref{ch:serving}).

\begin{finding}[two failures share one message --- still open]
An \emph{empty completion} (for example, a response blocked by the provider) returns the same ``I could not find anything about that\ldots'' text as an \emph{empty retrieval}. The reader is told the textbooks lack the topic when the real cause was the model call.
\end{finding}

\section{Citation validation}
\label{sec:citations}

``A language model asked to cite its sources will occasionally produce a plausible-looking marker for a passage it never saw --- and a fabricated citation is worse than none, because it is more convincing'' \ev{src/vidyarag/generate/citations.py}. The validator guarantees that every citation a reader sees points at a passage the system really supplied.

A marker is a bracket holding one- or two-digit numbers, separated by commas or semicolons. The first version accepted only \code{[1]}, but the model writes \code{[1, 2]} ``at least as often'', so grouped citations were silently dropped. The comment draws the lesson: ``Losing a real citation is the same class of failure as inventing one: the reference list stops matching the answer.''

\code{resolve_citations} maps each unique in-range marker to its passage. \code{strip_invalid_markers} removes out-of-range numbers from the text; within a group it removes only the invalid ones. Out-of-range markers are dropped, never remapped: ``\code{[7]} against five chunks is a hallucinated reference, and mapping it to something nearby would manufacture false provenance.''

\begin{measured}[constructed input, five passages supplied]
\code{...diffusion [2, 7] through GLUT proteins [3]. [9]} becomes \code{...diffusion [2] through GLUT proteins [3].}, with citations to passages 2 and 3 only. \code{[1][2]} yields two citations, \code{[ 4 ; 5 ]} is normalised to \code{[4, 5]}, and \code{[1-3]} or \code{[100]} are not markers, so they are left untouched.
\end{measured}

A valid marker proves only that a citation points at a supplied passage, not that the passage says what the sentence says. The test \code{test_every_citation_resolves_to_a_real_page_and_section} is, by its own docstring, ``the assertion the whole project's credibility rests on.''

\begin{finding}[the documented citation format is idealised --- still open]
The README and \code{ATTRIBUTION.md} show \texttt{Biology, \textsection 7.3, p.214}. The code renders \code{Biology, 7.3. <section title>*, p.214}: no section sign, the full section title, and the outline's asterisk.
\end{finding}
